% BHU PYQ -- Manually transcribed & error-corrected
\documentclass[11pt,a4paper]{article}

\usepackage[a4paper,top=2.5cm,bottom=2.5cm,left=2.8cm,right=2.8cm,headheight=14pt]{geometry}
\usepackage{amsmath,amssymb}
\usepackage[T1]{fontenc}
\usepackage[utf8]{inputenc}
\usepackage{lmodern}
\usepackage{microtype}
\usepackage{fancyhdr}
\usepackage{enumitem}
\usepackage{hyperref}
\usepackage{lastpage}
\usepackage{parskip}

\hypersetup{colorlinks=true,urlcolor=black,linkcolor=black,
  pdftitle={CHB-505: Environmental & Nuclear Chemistry -- BHU PYQ},pdfauthor={Banaras Hindu University}}

\pagestyle{fancy}\fancyhf{}
\renewcommand{\headrulewidth}{0.4pt}\renewcommand{\footrulewidth}{0.4pt}
\fancyhead[L]{\small   \textit{Banaras Hindu University}}
\fancyhead[R]{\small   \textit{CHB-505: Environmental \& Nuclear Chemistry}}
\fancyfoot[C]{\small Page  \thepage\ of \pageref{LastPage}}


\newlist{parts}{enumerate}{2}
\setlist[parts,1]{label=(\alph*),leftmargin=*,itemsep=5pt,topsep=3pt}
\setlist[parts,2]{label=(\roman*),leftmargin=*,itemsep=3pt,topsep=2pt}

\newcommand{\pts}[1]{\hfill   \textit{[#1]}}


\begin{document}


\begin{center}
  {\large\bfseries BANARAS HINDU UNIVERSITY}\\[2pt]
  {
\normalsize B.Sc. (Hons.) Semester V Examination, 2022-23}\\[6pt]
  \rule{0.8\linewidth}{0.5pt}\\[6pt]
  {\large\bfseries Paper No. CHB-505: Environmental \& Nuclear Chemistry}\\[4pt]
  {
\normalsize Time: 3 Hours\hfill Maximum Marks: 70}
\end{center}
\rule{\linewidth}{0.4pt}
\small



\noindent   \textit{Instructions: Answer five questions in total, including Question~1 from each section which is compulsory. Use separate answer books for Section~A and Section~B.}

\normalsize
\bigskip


\begin{center}
  {\large\bfseries SECTION A}\\[2pt]
  {
\normalsize (Environmental Chemistry -- Marks: 35)}
\end{center}
\rule{0.4\linewidth}{0.2pt}
\smallskip




\noindent   \textbf{Question 1.} Answer all parts of the following: \pts{5 $ \times$ 2 = 10}

\begin{parts}
  \item Differentiate between receptor and sink in environmental chemistry.
  \item Differentiate between lithosphere and biosphere.
  \item Discuss about the following: (i) Sludge, (ii) Van Allen Belts, and (iii) Pollution Pathways.
  \item Explain soot and contaminant.
  \item Discuss major regions of the atmosphere.
\end{parts}

\medskip\rule{\linewidth}{0.2pt}




\noindent   \textbf{Question 2.}

\begin{parts}
  \item Describe various environmental segments. How can we divide the major regions of the atmosphere? State temperature variations and chemical species present. \pts{8}
  \item Discuss the greenhouse effect and global warming. What are its major consequences? \pts{6}
\end{parts}

\medskip\rule{\linewidth}{0.2pt}




\noindent   \textbf{Question 3.} Account for any four of the following: \pts{4 $ \times$ 3.5 = 14}

\begin{parts}
  \item Minamata incident and Bhopal Gas Tragedy
  \item Wastewater treatment of domestic effluents
  \item Air pollution and important air quality parameters (BOD and COD)
  \item Auto-exhaust emissions and their catalytic converters controls
  \item Stratospheric ozone depletion chemistry
\end{parts}


\newpage

\begin{center}
  {\large\bfseries SECTION B}\\[2pt]
  {
\normalsize (Nuclear Chemistry -- Marks: 35)}
\end{center}
\rule{0.4\linewidth}{0.2pt}
\smallskip




\noindent   \textbf{Question 4.} Answer all parts of the following: \pts{5 $ \times$ 2 = 10}

\begin{parts}
  \item What is the basis of using radioisotopes for the age determination of historical and archaeological organic samples (radiocarbon dating)?
  \item What is the magic number of nucleons? Give two experimental evidences confirming the existence of magic numbers.
  \item Differentiate between nuclear reaction cross section and geometric cross section.
  \item Predict the modes of decay for neutron-rich and neutron-deficient nuclides.
  \item Explain the principle of isotopic dilution analysis.
\end{parts}

\medskip\rule{\linewidth}{0.2pt}




\noindent   \textbf{Question 5.}

\begin{parts}
  \item What are the different types of nuclear reactions? Describe radiative capture and transfer reactions in detail. \pts{6}
  \item Show that the kinetic energy available for exciting the target nucleus is the kinetic energy of the projectile and target in the Centre of Mass System (CMS). \pts{6}
  \item Write the precautions to be taken in the use of radioisotopes as tracers. \pts{2}
\end{parts}

\medskip\rule{\linewidth}{0.2pt}




\noindent   \textbf{Question 6.}

\begin{parts}
  \item Discuss the compound nucleus theory of nuclear reactions. Describe the Bohr model of nuclear fission. \pts{8}
  \item Discuss the nuclear shell model, showing magic numbers and spin-parity predictions. \pts{6}
\end{parts}

\vfill
\rule{\linewidth}{0.4pt}

\begin{center}\small
  \href{https://bhu-exam.vercel.app/}{Click here to visit our website}\\[4pt]
  \href{https://me-aryan.vercel.app/}{Click here to visit our contributor}\\[4pt]
  \href{https://quashan-me.vercel.app/}{Click here to visit our contributor}
\end{center}

\end{document}
